\exercisesection{Koszul Complexes}

\begin{enumerate}
\def\labelenumi{\arabic{enumi})}
\tightlist
\item
  Let A be a ring, L an A-module, and K the complex \(\mathbf{S}(\mathbf{L})\otimes_{\mathbf{A}}\symbf{\Lambda}(\mathbf{L})\) (X, p.~151, example 1).
\end{enumerate}

\begin{enumerate}
\def\labelenumi{\alph{enumi})}
\tightlist
\item
  Show that there exists an A-linear endomorphism s of K, graded of degree zero, such that
\end{enumerate}

\[
d s (x) + s d (x) = (p + q) x
\]

for all integers \(p, q \geqslant 0\) and every \(x \in \mathbf{S}^p L \otimes_{\mathbf{A}} \Lambda^q L\) .

\begin{enumerate}
\def\labelenumi{\alph{enumi})}
\setcounter{enumi}{1}
\tightlist
\item
  Assume that A is a Q-algebra. Show that K defines a resolution of the A-module A.
\end{enumerate}

\begin{enumerate}
\def\labelenumi{\arabic{enumi})}
\setcounter{enumi}{1}
\tightlist
\item
  Let A be a ring, \(u: \mathbf{L} \to \mathbf{M}\) a homomorphism of A-modules.
\end{enumerate}

\begin{enumerate}
\def\labelenumi{\alph{enumi})}
\item
  We denote by \(\varepsilon\) the commutation factor on \(\mathbf{Z}^2\) defined by \(\varepsilon(\alpha, \beta) = \alpha_2 \beta_2\) for \(\alpha, \beta \in \mathbf{Z}^2\) , \(\alpha = (\alpha_1, \alpha_2)\) , \(\beta = (\beta_1, \beta_2)\) . Show that there exists a unique (A, \(\varepsilon\) )-derivation (III, p.~118) \(d\) of the bigraded algebra \(\mathbf{S}(M) \otimes_A \symbf{\Lambda}(L)\) of degree \((1, -1)\) , such that one has \(d(1 \otimes x) = u(x) \otimes 1\) for \(x \in L\) . Show that \(d \circ d = 0\) , so that \(\mathbf{S}(M) \otimes_A \symbf{\Lambda}(L)\) , equipped with the grading induced by that of \(\symbf{\Lambda}(L)\) and with the differential \(d\) , defines a complex \(\mathbf{K}(u)\) .
\item
  If M = L and u = \(ld_{L}\) , show that \(\mathbf{K}(u)\) is equal to the complex defined in no. 3, p.~151.
\item
  Let \(\mathbf{B} = \mathbf{S}_{\mathbf{A}}(\mathbf{M})\) , and let \(\overline{u}: \mathbf{B} \otimes_{\mathbf{A}} \mathbf{L} \to \mathbf{B}\) be the B-homomorphism deduced from \(u\) . Show that \(\mathbf{K}(u)\) is canonically identified with the complex \(\mathbf{K}^{\mathbf{B}}(\overline{u})\) .
\item
  Let \(u': L' \to M'\) be a second A-homomorphism. Show that the complexes \(K(u \oplus u')\) and \(K(u) \otimes_{A} K(u')\) are isomorphic.
\item
  Let \(f: \mathbf{L} \to \mathbf{L}'\) and \(g: \mathbf{M} \to \mathbf{M}'\) be two A-homomorphisms such that \(g \circ u = u' \circ f\) . Show that the homomorphism \(\mathbf{S}(g) \otimes \symbf{\Lambda}(f)\) is a morphism of complexes from \(\mathbf{K}(u)\) to \(\mathbf{K}(u')\) . In particular, if \(v: \mathbf{M} \to \mathbf{P}\) is an A-homomorphism such that \(v \circ u = 0\) , we obtain a morphism of complexes \(h: \mathbf{K}(u) \to \mathbf{S}_{\mathbf{A}}(\mathbf{P})\) .
\item
  We assume that the sequence \(0 \rightarrow L \xrightarrow{u} M \xrightarrow{v} P \rightarrow 0\) is an exact sequence of flat A-modules. Show that h is an isomorphism (reduce to the case where L, M, P are finitely generated free, and use d)). Deduce from this for every \(n \geqslant 1\) an exact sequence
\end{enumerate}

\[
0 \rightarrow \Lambda^ {n} L \rightarrow M \otimes_ {A} \Lambda^ {n - 1} L \rightarrow \dots \rightarrow S ^ {n - 1} M \otimes_ {A} L \rightarrow S ^ {n} M \rightarrow S ^ {n} P \rightarrow 0.
\]

\begin{enumerate}
\def\labelenumi{\arabic{enumi})}
\setcounter{enumi}{2}
\tightlist
\item
  Let \(n\) be an integer, \(k\) a ring, \(A = k[X_1, \ldots, X_n]\) . Show that for every A-module M there exists an isomorphism of graded A-modules
\end{enumerate}

\[
\delta_ {\mathbf {M}}: \operatorname{Tor} ^ {\mathbf {A}} (k, \mathbf {M}) \rightarrow \operatorname{Ext} _ {\mathbf {A}} (k, \mathbf {M}) (- n)
\]

such that for every homomorphism of A-modules \(u: M \to N\) , one has Ext \((1_k, u) \circ \delta_M = \delta_N \circ \text{Tor } (1_k, u)\) . 4) Let A be a ring, L an A-module, \(u: L \to A\) a linear form.

\begin{enumerate}
\def\labelenumi{\alph{enumi})}
\item
  Show that the complex \(\mathbf{K}(u)\) , endowed with the algebra structure of \(\symbf{\Lambda}(\mathrm{L})\) , is a graded differential A-algebra (X, p.~183, exercise 18).
\item
  Let B be a differential graded A-algebra such that \(B_{0} = A\) and let \(f : L \to B_{1}\) be an A-homomorphism such that \(d_{1} \circ f = u\) . Show that there exists a unique morphism of graded differential A-algebras \(F : \mathbf{K}(u) \to \mathbf{B}\) such that \(F_{0} = 1_{A}\) and \(F_{1} = f\) .
\end{enumerate}

\begin{enumerate}
\def\labelenumi{\arabic{enumi})}
\setcounter{enumi}{4}
\item
  Let A be a ring, M an A-module, \(x = (x_1, \ldots, x_n)\) a sequence of elements of A, \(k = (k_1, \ldots, k_n) \in \mathbf{N}^n\) , \(x^k = (x_1^{k_1}, \ldots, x_{n}^{k_n})\) . Prove that the sequence \(x^k\) is M-regular if and only if the sequence \(x\) is; in that case the A-module \(\mathrm{M}/(x^k)\) admits a composition series whose quotients are isomorphic to \(\mathrm{M}/(x)\) . (Argue by induction on \(n\) .)
\item
  Let A be a ring, \(x = (x_1, \ldots, x_n)\) a sequence of elements of A, \(x\) the ideal \(\sum_{i} x_i\) A. Let
\end{enumerate}

\[
0 \rightarrow \mathbf {M} ^ {\prime} \stackrel {i} {\rightarrow} \mathbf {M} \rightarrow \mathbf {M} ^ {\prime \prime} \rightarrow 0
\]

an exact sequence of A-modules. Suppose that the sequence x is completely secant for M and that the A-modules \(\mathrm{M}''/(x_{1}\mathrm{M}''+\cdots+x_{i-1}\mathrm{M}'')\) are separated for the x-adic topology ( \(1\leqslant i\leqslant n\) ). Show that the following conditions are equivalent:

α) The sequence x is completely secant for M''.

β) The homomorphism \(\bigoplus_{r} (\mathfrak{x}^{r} \, \mathrm{M}^{\prime}/\mathfrak{x}^{r+1} \, \mathrm{M}') \to \bigoplus_{r} (\mathfrak{x}^{r} \, \mathrm{M}/\mathfrak{x}^{r+1} \, \mathrm{M})\) induced by \(i\) is injective.

\(\gamma)\) The homomorphism \(\mathbf{M}' / \mathfrak{x}\mathbf{M}' \to \mathbf{M} / \mathfrak{x}\mathbf{M}\) induced by \(i\) is injective.

\begin{enumerate}
\def\labelenumi{\arabic{enumi})}
\setcounter{enumi}{6}
\tightlist
\item
  Let A be a ring, M an A-module, \(x = (x_1, \ldots, x_n)\) a sequence of elements of A.
\end{enumerate}

\[
(X, p p. 1 7 5 - 1 7 7,
\]

\[
\mathrm{E} _ {p q} ^ {2} = \operatorname{Tor} _ {p} ^ {\mathrm{A}} \left(\mathrm{H} _ {q} (\boldsymbol {x}, \mathrm{A}) \right.
\]

\begin{enumerate}
\def\labelenumi{\alph{enumi})}
\setcounter{enumi}{1}
\tightlist
\item
  Let \(p\) be an integer \(\leqslant n\) ; set \(x' = (x_1, \ldots, x_p)\) and \(x'' = (x_{p+1}, \ldots, x_n)\) . Show that there exists a spectral sequence E such that 'E \(_{pa}^{2}\) = H \(_{p}\) x'`, H \(_{a}\) (x', M)), converging to H. (x, M).
\end{enumerate}

\begin{enumerate}
\def\labelenumi{\arabic{enumi})}
\setcounter{enumi}{7}
\tightlist
\item
  Let A be a ring, M an A-module, \(\boldsymbol{x}' = (x_1, \ldots, x_p)\) and \(\boldsymbol{x}'' = (x_{p+1}, \ldots, x_n)\) two sequences of elements of A; we denote by \(\boldsymbol{x}\) the sequence \((x_1, \ldots, x_n)\) , and by \(\boldsymbol{M}'\) the A-module \(\mathrm{M}/(\boldsymbol{x}')\) M.
\end{enumerate}

\begin{enumerate}
\def\labelenumi{\alph{enumi})}
\item
  If \(x'\) is completely secant for M and \(x''\) completely secant for M', show that x is completely secant for M.
\item
  If the sequences \(x\) and \(x'\) are completely secant for M, show that the sequence \(x''\) is completely secant for M'.
\item
  Give an example of an A-regular sequence \((x, y)\) contained in the radical of A, such that x is completely secant for A/yA, but y is not completely secant for A.
\end{enumerate}

\begin{enumerate}
\def\labelenumi{\arabic{enumi})}
\setcounter{enumi}{8}
\tightlist
\item
  Let A be a ring, M an A-module, and k an integer \(\geqslant\) 1. One says that a family x of elements of A is k-secant for M if one has \(H_{i}(x, M) = 0\) for \(1 \leqslant i \leqslant k\) . A family x is therefore completely secant if and only if it is k-secant for every \(k \geqslant 1\) .
\end{enumerate}

\begin{enumerate}
\def\labelenumi{\alph{enumi})}
\item
  Let \(0 \to M' \to M \to M'' \to 0\) be an exact sequence of A-modules. Show that if the family \(x\) is \(k\) -secant for \(M'\) and \(M''\) , it is so for \(M\) .
\item
  If x is k-secant for M and if N is a flat A-module, then x is k-secant for M \(\otimes_{A} N\) .
\item
  Let \(a_{1},\ldots,a_{n}\) be integers \(\geqslant 1\) . Show that a sequence \((x_{1},\ldots,x_{n})\) is \(k\) -secant for M if and only if the same holds for the sequence \((x_{1}^{a_{1}},\ldots,x_{n}^{a_{n}})\) .
\item
  For every pair of integers k, n with \(1 \leqslant k < n\) , give an example of a ring A, an A-module M, and a sequence \((x_{1}, \ldots, x_{n})\) of elements of A that is k-secant but not \((k + 1)\) -secant for M (use X, p.~159, remark 4).
\end{enumerate}

\begin{enumerate}
\def\labelenumi{\arabic{enumi})}
\setcounter{enumi}{9}
\tightlist
\item
  Let A be a ring, M an A-module, \(\boldsymbol{x} = (x_{1}, \ldots, x_{n})\) a sequence of elements of A, p an integer \(\leqslant n\) . We write \(\boldsymbol{x}' = (x_{1}, \ldots, x_{p})\) , \(\boldsymbol{x}'' = (x_{p+1}, \ldots, x_{n})\) and \(\mathrm{M}' = \mathrm{M}/(\boldsymbol{x}')\) M.
\end{enumerate}

\begin{enumerate}
\def\labelenumi{\alph{enumi})}
\item
  Let \(k, k'\) be two integers such that \(k \leqslant k'\) . Suppose that the sequence \(x'\) is \(k'\) -secant for M (exercise 9). Show that the sequence \(x\) is \(k\) -secant for M if and only if the sequence \(x''\) is \(k\) -secant for M'.
\item
  Show that if the sequence x is 1-secant for M, then the sequence \(x''\) is 1-secant for \(M'\) .
\end{enumerate}

¶ 11) Let A be a noetherian local ring, m its maximal ideal, k = A/m. Let M be a nonzero finitely generated A-module. The depth of M, denoted prof (M), is defined as the supremum of the integers n for which there exists a sequence \((x_{1}, \ldots, x_{n})\) of elements of m that is completely secant for M.

\begin{enumerate}
\def\labelenumi{\alph{enumi})}
\item
  Show that prof (M) = 0 if and only if m is associated to M (X, p.~172, exercise 28), in other words, if M contains a submodule isomorphic to k (observe that an ideal contained in a union of prime ideals is necessarily contained in one of them).
\item
  Show that prof(M) is the infimum of the integers n for which Ext \(_{A}^{n}\) (k, M) is nonzero (argue by induction on n). If prof(M) = n, then every sequence (x \(_{1}\) , \ldots, x \(_{p}\) ) of elements of m that is completely secant for M (p ≤ n) can be extended to a sequence (x \(_{1}\) , \ldots, x \(_{n}\) ) in m completely secant for M. c) If (x \(_{1}\) , \ldots, x \(_{k}\) ) is a sequence of elements of m completely secant for M, show that
\end{enumerate}

\[
\operatorname{prof} \left(\mathbf {M} / (x _ {1} \mathbf {M} + \dots + x _ {k} \mathbf {M})\right) = \operatorname{prof} (\mathbf {M}) - k.
\]

\begin{enumerate}
\def\labelenumi{\alph{enumi})}
\setcounter{enumi}{3}
\item
  Assume that M has finite projective dimension. Prove the equality dpA(M) + prof(M) = prof(A) (argue by induction on dpA(M)).
\item
  Show that prof(A) is the supremum of the projective dimensions of A-modules that are finitely generated and of finite projective dimension.
\end{enumerate}

¶ 12) Let A be a noetherian local ring, m its maximal ideal.

\begin{enumerate}
\def\labelenumi{\alph{enumi})}
\tightlist
\item
  Show that if \(0 < \mathrm{dh}(\mathbf{A}) < \infty\) there exists a non-zero-divisor element in \(\mathfrak{m} - \mathfrak{m}^2\) (show as in Exercise 11, a) that \(\mathfrak{m}\) would otherwise be associated with \(\mathbf{A}\) whence an exact sequence
\end{enumerate}

\[
0 \to k \to \mathbf {A} \to \mathbf {N} \to 0;
\]

lead to a contradiction by using Exercise 16, p.~203).

\begin{enumerate}
\def\labelenumi{\alph{enumi})}
\setcounter{enumi}{1}
\tightlist
\item
  For any integer n, show that dh(A) = n if and only if there exists a sequence \((x_{1}, \ldots, x_{n})\) completely secant for A, generating m (argue by induction on n, using a) and exercise 14, p.~203).
\end{enumerate}

A local ring is said to be regular if it is noetherian and has finite homological dimension.

\begin{enumerate}
\def\labelenumi{\arabic{enumi})}
\setcounter{enumi}{12}
\tightlist
\item
  Let A be a regular local ring (Exercise 14) of homological dimension n; denote its maximal ideal by m and let k = A/m.
\end{enumerate}

\begin{enumerate}
\def\labelenumi{\alph{enumi})}
\item
  Show that one has \(\dim_k(m/m^2) = n\) .
\item
  Show that the \(k\) -algebra \(\bigoplus_{r\geqslant 0}(m^{r} / m^{r + 1})\) is isomorphic to a polynomial algebra in \(n\) indeterminates.
\item
  If \(n = 0\) , A is a field. If \(n = 1\) , show that A is a principal ideal ring * and therefore a discrete valuation ring * (show that A is an integral domain by observing that \(m \cap m^n = \bigcap m^n\) , whence \(\bigcap m^n = 0\) ).
\item
  Show that the ring A{[}{[}T{]}{]} is regular.
\end{enumerate}

\begin{enumerate}
\def\labelenumi{\arabic{enumi})}
\setcounter{enumi}{13}
\item
  Let A be a ring, and let a be an ideal of A. Consider the graded alternating (A/a)-algebra \(\mathrm{Tor}^{\mathbf{A}}(\mathbf{A}/\mathbf{a}, \mathbf{A}/\mathbf{a})\) (cf. X, p. 201, exercise 9) and the homomorphism of graded (A/a)-algebras \(\theta: \Lambda_{\mathbf{A}/\mathbf{a}}(\mathbf{a}/\mathbf{a}^{2}) \to \mathrm{Tor}^{\mathbf{A}}(\mathbf{A}/\mathbf{a}, \mathbf{A}/\mathbf{a})\) deduced from the isomorphism \(a/a^{2} \to \mathrm{Tor}_{1}^{\mathbf{A}}(\mathbf{A}/\mathbf{a}, \mathbf{A}/\mathbf{a})\) (X, p. 73). Show that if the ideal a is generated by a completely secant sequence in A, \(\theta\) is an isomorphism.
\item
  Let A be a noetherian ring, and let a be an ideal contained in the radical of A. Suppose that the (A/a)-module \(a/a^{2}\) is free and that the homomorphism \(\theta_{2}: \Lambda_{A/a}^{2}(a/a^{2}) \to \operatorname{Tor}_{2}^{\mathbf{A}}(\mathbf{A}/\mathbf{a}, \mathbf{A}/\mathbf{a})\) (exercise 14) is surjective. Show that a is generated by a completely secant sequence in A. (If x is a minimal family of generators of a, show with the aid of exercise 7 (or directly) that Coker ( \(\theta_{2}\) ) is isomorphic to
\end{enumerate}

\[
\mathrm{H} _ {1} (\boldsymbol {x}, \mathbf {A}) \otimes_ {\mathbf {A}} \mathbf {A} / \mathfrak {a}.
\]

\begin{enumerate}
\def\labelenumi{\arabic{enumi})}
\setcounter{enumi}{15}
\tightlist
\item
  Let A be a local noetherian ring, m its maximal ideal, k = A/m. Show that the homomorphism of graded k-algebras \(\theta : \symbf{\Lambda}_{k}(\mathfrak{m}/\mathfrak{m}^{2}) \to \operatorname{Tor}^{\mathbf{A}}(k, k)\) is injective: if \(\theta_{2}\) is bijective, the ring A is regular (exercise 12).
\end{enumerate}

(If x is a minimal system of generators of m, show that the complex K(x, A) satisfies the conditions of exercise 14, p. 182; use exercises 4 and 15.)

¶ 17) Let A be a ring such that the A-module \(A_{s}\) has finite injective dimension (X, p. 202, exercise 4; a regular local ring (exercise 12) satisfies this property).

\begin{enumerate}
\def\labelenumi{\alph{enumi})}
\tightlist
\item
  Let x be a non-zero-divisor and non-invertible element of A. Show that one has
\end{enumerate}

\[
\mathrm{di} _ {\mathbf {A} / \mathbf {A} x} (\mathbf {A} / \mathbf {A} x) = \mathrm{di} _ {\mathbf {A}} (\mathbf {A}) - 1.
\]

(Let 0 → A → I⁰ → I¹ → \ldots{} be a minimal injective resolution (X, p. 182, exercise 13) of A; show that the complex Homₐ(A/AX, I¹) → Homₐ(A/AX, I²) → \ldots{} defines a minimal injective resolution of the (A/AX)-module A/AX.)

\begin{enumerate}
\def\labelenumi{\alph{enumi})}
\setcounter{enumi}{1}
\item
  We henceforth suppose that the ring A is local noetherian. Prove the equality di(A) = prof(A) (exercise 11; reason by induction on di(A)).
\item
  Let \((x_{1},\ldots ,x_{n})\) be a completely secant sequence for A, contained in the maximal ideal of A, with \(n = \mathrm{di}(\mathbf{A})\) . Show that the ring A is self-injective (X, p. 171, exercise 26).
\end{enumerate}

\begin{enumerate}
\def\labelenumi{\arabic{enumi})}
\setcounter{enumi}{17}
\tightlist
\item
  Let A be a noetherian ring, C a stable set of classes of A-modules (X, p. 40), \(x = (x_{1}, \ldots, x_{n})\)
\end{enumerate}

a sequence of elements of A, x the ideal generated by x. We denote by \(\mathcal{C}_{x}\) the set of elements of \(\mathcal{C}\) which are annihilated by x. Let M be a finitely generated A-module.

\begin{enumerate}
\def\labelenumi{\alph{enumi})}
\tightlist
\item
  Prove that for every \(i \geqslant 0\) , there exists an element \(\mathbf{M}_{i}\) of \(\mathcal{C}_{x}\) such that one has
\end{enumerate}

\[
\sum_ {k \geqslant 0} (- 1) ^ {k} \left[ \mathrm{H} _ {i + k} (\mathbf {x}, \mathrm{M}) \right] = [ \mathrm{M} _ {i} ]
\]

in \(\mathbf{K}(\mathcal{C}_{\pmb{x}})\) (reason by induction on \(n\) , using prop. 4, p. 157; in the case \(n = 1\) , consider the successive powers of the endomorphism \((x_{1})_{\mathbb{M}}\) ).

\begin{enumerate}
\def\labelenumi{\alph{enumi})}
\setcounter{enumi}{1}
\tightlist
\item
  We assume that \(\mathbf{A} / x\) has finite length. Show that there exist positive integers \(m_{i}\) ( \(i \geqslant 0\) ) such that one has length \(\mathrm{H}_{i}(x, \mathbf{M}) = m_{i} + m_{i+1}\) .
\end{enumerate}

